%============================================================
% UKURAN FONT
% Semua ukuran font template ada di sini. Ubah satu baris,
% berlaku di semua halaman yang memakai peran tersebut.
%
% Format: \fontsize{ukuran huruf}{jarak baris}\selectfont (pt)
% Nilai di bawah adalah ukuran Buku PA. Proposal & Progres
% meng-override sebagian lewat \renewcommand di
% formats/format_proposal.tex dan formats/format_progres.tex.
%
% Di-input setelah config/packages.tex dan sebelum
% config/formatting.tex di setiap main_*.tex.
%============================================================

% ── Judul Isi Dokumen (Buku PA; Proposal/Progres: 12pt) ──────
\newcommand{\FontJudulBab}{\fontsize{14}{17}\selectfont}           % BAB n + judul bab
\newcommand{\FontJudulSubbab}{\fontsize{14}{17}\selectfont}        % 1.1 (section)
\newcommand{\FontJudulSubSubbab}{\fontsize{12}{14}\selectfont}     % 1.1.1 (subsection)
\newcommand{\FontJudulSubSubSubbab}{\fontsize{12}{14}\selectfont}  % tanpa nomor (subsubsection)
\newcommand{\FontJudulDaftar}{\fontsize{14}{17}\selectfont}        % DAFTAR ISI / GAMBAR / TABEL
\newcommand{\FontKutipanPanjang}{\fontsize{10}{12}\selectfont}     % lingkungan kutipanpanjang
\newcommand{\FontSumber}{\fontsize{12}{15}\selectfont}             % "Sumber: ..." di bawah gambar/tabel

% ── Cover Buku PA (formats/cover_buku.tex) ───────────────────
\newcommand{\FontCoverBukuLabel}{\fontsize{20}{24}\selectfont}     % "PROYEK AKHIR"
\newcommand{\FontCoverBukuJudul}{\fontsize{18}{20}\selectfont}     % judul PA
\newcommand{\FontCoverBukuIsi}{\fontsize{16}{20}\selectfont}       % nama, NRP, pembimbing, prodi

% ── Halaman Judul Dalam Buku PA (formats/halaman_judul.tex) ──
\newcommand{\FontJudulDalamLabel}{\fontsize{20}{0}\selectfont}     % "PROYEK AKHIR"
\newcommand{\FontJudulDalamJudul}{\fontsize{18}{20}\selectfont}    % judul PA
\newcommand{\FontJudulDalamIsi}{\fontsize{16}{20}\selectfont}      % nama, NRP, pembimbing, prodi

% ── Halaman Pengesahan (formats/pengesahan.tex) ──────────────
\newcommand{\FontPengesahanLabel}{\fontsize{20}{0}\selectfont}     % "HALAMAN PENGESAHAN"
\newcommand{\FontPengesahanJudul}{\fontsize{18}{16}\selectfont}    % judul PA
\newcommand{\FontPengesahanIsi}{\fontsize{14}{17}\selectfont}      % isi & tabel tanda tangan

% ── Halaman Pernyataan (orisinalitas, hak cipta) ─────────────
\newcommand{\FontPernyataanLabel}{\fontsize{20}{10}\selectfont}    % judul halaman pernyataan
\newcommand{\FontPernyataanIsi}{\fontsize{12}{15}\selectfont}      % isi pernyataan orisinalitas

% ── Cover Proposal & Progres (formats/cover_proposal|progres.tex) ──
\newcommand{\FontCoverRingkasLabel}{\fontsize{20}{24}\selectfont}       % "PROPOSAL/PROGRES PROYEK AKHIR"
\newcommand{\FontCoverRingkasIsi}{\fontsize{18}{22}\selectfont}         % judul, nama, NRP, prodi
\newcommand{\FontCoverRingkasPembimbing}{\fontsize{18}{20}\selectfont}  % nama & NIP pembimbing

% Catatan: label 20pt memakai jarak baris berbeda (24pt di cover,
% 0pt di halaman judul & pengesahan, 10pt di halaman pernyataan).
% Nilai ini dipertahankan apa adanya agar hasil PDF tidak berubah.
% Jarak baris hanya berpengaruh jika judul tersebut lebih dari 1 baris.
